%% ma: L11-B14-0011
%% lop: 11
%% chuong: 4
%% bai: 14
%% dang: 11.14.1
%% loai: TN4
%% muc: TH
%% dap_an: "A"
%% nguon: "tai-lieu-hoc-tap-mon-toan-lop-11-theo-chuong-trinh-giao-duc-pho-thong-2018.pdf (Lớp Toán Thầy Duy, Phạm Lê Duy, THPT Chu Văn An), Chương 4 Bài 5, Phần 1, Câu 6 (tr.526)"
%% kien_thuc: "Lớp 11 Bài 14 (phép chiếu song song)"
%% trang_thai: chinh-thuc
%% ghi_chu: "Đã duyệt 03/11/2026."
\begin{ex}
Cho tứ diện $ABCD$. Gọi $G$ là trọng tâm của tam giác $ACD$. Hình chiếu $G^{\prime}$ của điểm $G$ trên mặt phẳng $(BCD)$ theo phương chiếu $AB$ là
\choice{\True Trọng tâm của tam giác $BCD$.}{Trung điểm $BC$.}{Trung điểm $CD$.}{Trung điểm $BI$ với $I$ là trung điểm $CD$.}
\loigiai{\sloppy Theo phương $AB$ lên $(BCD)$: $A\mapsto B$, $C\mapsto C$, $D\mapsto D$. Phép chiếu song song bảo toàn tỉ số độ dài của các đoạn thẳng trên cùng một đường thẳng, nên trọng tâm $G$ của $\triangle ACD$ (giao các trung tuyến) có hình chiếu là trọng tâm của $\triangle BCD$.}
\end{ex}
